\phantomsection\label{app:lower-coefficients-proof}

\begin{proof}[Proof of Lemma~\ref{p4:transverse-coefficients}]
For $A=1$ the whole $C>1$ family has the invariant parabola, so
$\Delta_0(1,C)=0$.
Write $K=c-1$ and
\[
 \mathcal E_c(X)=((X+1)^2+K)^2
 \exp\left[-\frac4{\sqrt K}
                \arctan\frac{X+1}{\sqrt K}\right].
\]
For $W=Y-(X^2-c)/2$, the exact affine equation at $\mu=0,C=c$ is
\[
 \dot W=2XW+(1-A)X^2,\qquad
 \mathcal E_c'/\mathcal E_c=2X/f_c.
\]
The graph variation in the $A$ direction therefore satisfies
\[
 W_{A,X}=(2X/f_c)W_A-X^2/f_c.
\]
The tail bound $z_A=O(q^3)$ gives $W_A=O(X)$ at infinity.
Since $\mathcal E_c\asymp X^4$, the free homogeneous coefficient
vanishes, and the selected variations are
\[
 W_A^-(X)=-\mathcal E_c(X)\int_{-\infty}^X
       \frac{t^2}{f_c(t)\mathcal E_c(t)}\dd t,\qquad
 W_A^+(X)=\mathcal E_c(X)\int_X^\infty
       \frac{t^2}{f_c(t)\mathcal E_c(t)}\dd t.
\]
The positive integrand is $O(|t|^{-4})$ at both ends, so
$W_A^+-W_A^-$ is positive at the positive port. At fixed $v=d$,
\[
 W_z=\frac{d^2-z_c(d)}{z_c(d)^3}>0,
\]
conversion to raw height gives $\partial_A\Delta_0>0$.

Truncate the limiting tails at $v=\pm a$. The outer derivative is
$E_0'$ times the negative-tail first variation divided by the positive
one, tending to $\Lambda_0=E_0'\ell_-/\ell_+$ as $a\downarrow0$.
The parabola endpoints are $X=\pm X_a$, $X_a=a/z_c(a)\to\infty$,
with equal fixed-$v$ conversion factors $W_z$. Thus the derivative
is $\mathcal E_c(X_a)/\mathcal E_c(-X_a)$: its polynomial ratio tends
to one and its arctangent difference to $\pi$, giving
$\Lambda_0(1,c)=e^{-4\pi/\sqrt K}$. For every $C$ in the rectangle,
$\Lambda_0(1,C)=e^{-4\pi/\sqrt{C-1}}$, so differentiation gives
$\partial_C\Lambda_0(1,c)=2\pi K^{-3/2}e^{-4\pi/\sqrt K}$.
\end{proof}
